\BourbakiExercises{Exercises}

\begin{enumerate}
\def\labelenumi{\arabic{enumi})}
\tightlist
\item
  Let B be an algebra of finite degree over the commutative field K, and let b be an element of B.
\end{enumerate}

\begin{enumerate}
\def\labelenumi{\alph{enumi})}
\item
  If \(b\) is nilpotent, then we have \(\mathrm{Tr}_{\mathrm{M / K}}(b) = 0\) for every B-module M of finite dimension over K. Deduce that if \(b\) belongs to the radical of B, then we have \(\mathrm{Tr}_{\mathrm{M / K}}(bc) = 0\) for every \(c \in \mathrm{B}\) .
\item
  Prove that if \(\operatorname{Tr}_{\mathrm{S}/\mathrm{K}}(bc)=0\) for every simple B-module S and every \(c\in B\) , then b belongs to the radical of B (use the density theorem).
\end{enumerate}

\begin{enumerate}
\def\labelenumi{\arabic{enumi})}
\setcounter{enumi}{1}
\tightlist
\item
  Let B be an algebra of finite degree n over the commutative field K. For any element b of B, denote by \(\mathrm{Pm}_{\mathrm{B}/\mathrm{K}}(b;\mathrm{X})\) the minimal polynomial of b over K (V, §3, No.~1, p.~16). We have \(\mathrm{Pm}_{\mathrm{B}/\mathrm{K}}(b;\mathrm{X})=\mathrm{Pm}_{\mathrm{B}^{\circ}/\mathrm{K}}(b;\mathrm{X})\) .
\end{enumerate}

\begin{enumerate}
\def\labelenumi{\alph{enumi})}
\item
  Prove that \(\mathrm{Pm}_{\mathrm{B / K}}(b;\mathrm{X})\) divides \(\mathrm{Pc}_{\mathrm{B / K}}(b;\mathrm{X})\) , which divides \((\mathrm{Pm}_{\mathrm{B / K}}(b;\mathrm{X}))^n\) .
\item
  For every extension L of K, we have \(\mathrm{Pm}_{\mathrm{B}_{(\mathrm{L})}/\mathrm{L}}(1 \otimes b; \mathrm{X}) = \mathrm{Pm}_{\mathrm{B}/\mathrm{K}}(b; \mathrm{X})\) .
\item
  Let \((e_i)_{i\in \mathrm{I}}\) be a basis of B over K and \(\mathbf{Y} = (\mathrm{Y}_i)_{i\in \mathrm{I}}\) a family of variables. The principal polynomial of the algebra B with respect to the basis \((e_i)\) , denoted by \(\mathrm{P}((e_i);\mathrm{X};\mathbf{Y})\) or simply \(\mathrm{P}(\mathrm{X};\mathbf{Y})\) , is the minimal polynomial over \(\mathrm{K}(\mathbf{Y})\) of the element \(\sum_{i\in \mathrm{I}}\mathrm{Y}_i e_i\) of \(\mathrm{B}_{(\mathrm{K}(\mathbf{Y}))}\) . *Prove that it belongs to \(\mathrm{K}[\mathrm{X};\mathbf{Y}]\) (cf.~AC, VII, §3, n° 5, p.~230, lemma 1 and théorème 2).
\item
  Let \((e_{i}^{\prime})\) be another basis of B over K and \((\lambda_{ij})\) the matrix with entries in K such that \(e_{i} = \sum_{j} \lambda_{ij} e_{j}^{\prime}\) for \(i \in I\) ; if we set \(Y_{i}^{\prime} = \sum_{j} \lambda_{ji} Y_{j}\) , then \(\mathrm{P}((e_{i}); \mathrm{X}; \mathbf{Y}) = \mathrm{P}((e_{i}^{\prime}); \mathrm{X}; \mathbf{Y}^{\prime})\) . Let \(b = \sum_{i} \beta_{i} e_{i}\) be an element of B; the polynomial \(\mathrm{P}((e_{i}); \mathrm{X}; \beta_{1}, \ldots, \beta_{n})\) is independent of the choice of the basis \((e_{i})\) ; we call it the principal polynomial of b over K. It divides \(\mathrm{Pc}_{\mathrm{B/K}}(b; \mathrm{X})\) and is a multiple of \(\mathrm{Pm}_{\mathrm{B/K}}(b; \mathrm{X}).*\)
\item
  Let \(B'\) be another K-algebra of finite degree; calculate the principal polynomial of the algebra \(B \times B'\) for a suitable basis.
\end{enumerate}

\begin{enumerate}
\def\labelenumi{\arabic{enumi})}
\setcounter{enumi}{2}
\tightlist
\item
  \begin{enumerate}
  \def\labelenumii{\alph{enumii})}
  \tightlist
  \item
    Let A be a central simple algebra of finite degree over K. Prove that the principal polynomial of an element \(a\) of A is \(\mathrm{Pcrd}_{\mathrm{A} / \mathrm{K}}(a; \mathrm{X})\) . The principal polynomial of A (with respect to a basis of A over K) is the polynomial P considered in Lemma 4 of VIII, p.~345 (reduce to the case \(\mathrm{B} = \mathbf{M}_n(\mathrm{K})\) ).
  \end{enumerate}
\end{enumerate}

\begin{enumerate}
\def\labelenumi{\alph{enumi})}
\setcounter{enumi}{1}
\tightlist
\item
  Let B be an absolutely semisimple algebra over K and \(\mathcal{S}\) be the set of classes of simple B-modules. Prove that the principal polynomial of an element b of B is equal to the product of the polynomials \(\mathrm{Pc}_{\mathrm{S}/\mathrm{K}}(b;\mathrm{X})\) for S running through \(\mathcal{S}\) (reduce to the case when K is algebraically closed).
\end{enumerate}

¶ 4) Let D be a field, \(\sigma\) an automorphism of D, A the ring D{[}X{]} \(_{\sigma}\) (VIII, p.~11), \(n\) an integer \(\geqslant 2\) , and \(a\) a nonzero element of D. Denote the element X \(^{n}\) - \(a\) of A by \(f\) . \(a)\) The ideal A \(f\) is two-sided if and only if we have \(\sigma(a) = a\) and \(\sigma^{n}(x) = axa^{-1}\) for every \(x \in \mathrm{D}\) .

\begin{enumerate}
\def\labelenumi{\alph{enumi})}
\setcounter{enumi}{1}
\tightlist
\item
  Suppose that the group \(\mu_{n}\) of n-th roots of unity contained in the center of D is cyclic of order n and fixed by \(\sigma\) . Let \(g \in A\) be an irreducible monic polynomial (VIII, p.~22, Exercise 27) dividing f on the right. Prove that the degree d of g divides n and that f is the product of n/d polynomials of degree d.~(Let h be the monic polynomial such that \(Ah = \bigcap_{\zeta \in \mu_{n}} Ag_{\zeta}\) , where \(g_{\zeta}(X) = g(\zeta X)\) . Observe that we necessarily have \(h(\zeta X) = h(X)\) for every \(\zeta \in \mu_{n}\) , and deduce that h = f.~Then, apply loc. cit., b) by constructing, for \(1 \leqslant i \leqslant n/d\) , an element \(\zeta_{i}\) of \(\mu_{n}\) and a polynomial \(h_{i}\) of degree id such that \(Ag_{\zeta_{1}} \cap \cdots \cap Ag_{\zeta_{i}} = Ah_{i}\) .
\end{enumerate}

If, moreover, the ideal Af is two-sided, then the ring A/Af is semisimple, and its simple components all have the same length (observe that every simple \((\mathrm{A}/\mathrm{Af})\) -module is isomorphic to one of the modules \(A/Ag_{\zeta}\) ).

\begin{enumerate}
\def\labelenumi{\alph{enumi})}
\setcounter{enumi}{2}
\tightlist
\item
  Assume from now on that the hypotheses of b) are satisfied, that the ideal A f is two-sided, and that n is a prime number. Deduce from b) that we are in one of the following situations:
\end{enumerate}

\begin{enumerate}
\def\labelenumi{(\roman{enumi})}
\item
  There exists a \(b \in D\) such that \(a = b^{n}\) and that \(\sigma\) is the inner automorphism associated with b. The ring C = A/Af is isomorphic to the product ring \(D^{n}\) .
\item
  The condition of (i) is not satisfied, but there exists a \(c \in \mathrm{D}\) such that \(a = \sigma^{n-1}(c)\sigma^{n-2}(c)\ldots\sigma(c)c\) . The ring C is simple of length \(n\) .
\item
  There is no element c of D such that \(a = \sigma^{n-1}(c) \sigma^{n-2}(c) \cdots \sigma(c) c\) . The ring C is a field.
\end{enumerate}

In the case when D is commutative, recover the results of Exercise 11 of VIII, p.~330. If, moreover, n = 2, then C is a quaternion algebra over the subfield K of D invariant under \(\sigma\) .

\begin{enumerate}
\def\labelenumi{\alph{enumi})}
\setcounter{enumi}{3}
\item
  Prove that there exists an automorphism \(\tau\) of C of order \(n\) whose subfield of invariants is equal to D. In cases (ii) and (iii) above, D is a Galois subfield of C (Exercise 16 of VIII, p.~274).
\item
  Suppose that \(\sigma\) is the inner automorphism associated with an element \(d\) of D. Prove that we have \(a = d^n z\) , where \(z\) belongs to the center Z of C, and that the ring C is isomorphic to the tensor product \(\mathrm{D} \otimes_{\mathbb{Z}} \mathrm{Z}[\mathrm{X}] / (\mathrm{X}^n - z)\) .
\item
  Suppose that \(\sigma\) is not an inner automorphism. Prove that the center of C is the subfield of Z consisting of the elements invariant under \(\sigma\) .
\end{enumerate}

¶ 5) Let \(K_{0}\) be the field of rational fractions \(\mathbf{Q}(\mathrm{X})\) . Let K be the quadratic extension \(\mathrm{K}_{0}(\rho)\) , where \(\rho\) is a primitive cubic root of unity, D be the cyclic extension \(\mathrm{K}(\sqrt[3]{\mathrm{X}})\) of K, and \(\sigma\) be a generator of the Galois group of D over K. Consider the algebra C defined as in Exercise 4, with n = 3 and a = 2.

\begin{enumerate}
\def\labelenumi{\alph{enumi})}
\item
  Prove that C is a field with center K, of degree 9 over K. (To see that D has no element of norm 2, reduce to proving that there cannot be any relations of the form \(2F^{3} = P^{3} + XQ^{3} + X^{2}R^{3} + 3XPQR\) , where F, P, Q, R are mutually prime polynomials in \(\mathbf{Q}(\rho)[\mathrm{X}]\) . To do this, consider the lowest-degree terms of these polynomials.)
\item
  Prove that there does not exist any automorphism of C extending the unique \(K_{0}\) -automorphism of K distinct from the identity (reduce to the case when such an automorphism preserves D). Deduce that \(K_{0}\) is not a Galois subfield of C (VIII, p.~274, Exercise 16) and that there does not exist a \(K_{0}\) -automorphism of C of order 2.
\item
  Let \(\xi\) be an element of D such that \(\xi^3 = \mathrm{X}\) , and let \(\eta\) be an element of C such that \(\eta^3 = 2\) and \(\eta \xi = \rho \xi \eta\) ; set \(\zeta = \xi + \eta + \eta^2\) . Prove that \(\mathrm{K}(\zeta)\) is a maximal commutative subfield of D and is not Galois over K, whereas the maximal commutative subfields \(\mathrm{K}(\xi)\) and \(\mathrm{K}(\eta)\) are Galois over K (calculate \(\zeta^3\) ).
\end{enumerate}

\begin{enumerate}
\def\labelenumi{\arabic{enumi})}
\setcounter{enumi}{5}
\tightlist
\item
  Let D be a field with center K of degree \(m^{2}\) over K. Let x be an element of D of degree n \textgreater{} 1 over K, and let \(f(\mathrm{X})\) be its minimal polynomial over K.
\end{enumerate}

\begin{enumerate}
\def\labelenumi{\alph{enumi})}
\item
  Prove that a monic polynomial in D{[}X{]} (VIII, p.~11) of degree \textless{} n cannot be right divisible by all polynomials \(X - txt^{-1}\) for \(t \in D^{*}\) (observe that the monic generator of the ideal \(\cap(\mathrm{D}[\mathrm{X}](\mathrm{X}-txt^{-1}))\) has coefficients in K).
\item
  Let \(u(\mathrm{X}), v(\mathrm{X}), w(\mathrm{X})\) be three monic polynomials in D{[}X{]} such that \(u(\mathrm{X}) = v(\mathrm{X})w(\mathrm{X})\) and \(\mathrm{X} - x\) is a right divisor of \(u(\mathrm{X})\) but not of \(w(\mathrm{X})\) . Write \(w(\mathrm{X}) = q(\mathrm{X})(\mathrm{X} - x) + t\) with \(t \in \mathrm{D}^*\) . Prove that \(\mathrm{X} - txt^{-1}\) is a right divisor of \(v(\mathrm{X})\) .
\item
  Deduce from a) and b) that there exist n elements \(t_{1},\ldots,t_{n}\) of D such that we have
\end{enumerate}

\[
f (\mathrm{X}) = \left(\mathrm{X} - t _ {1} x t _ {1} ^ {- 1}\right) \dots \left(\mathrm{X} - t _ {n} x t _ {n} ^ {- 1}\right)
\]

(observe that \(f(\mathrm{X})\) is right divisible by all polynomial \(\mathrm{X} - txt^{-1}\) ).

\begin{enumerate}
\def\labelenumi{\alph{enumi})}
\setcounter{enumi}{3}
\tightlist
\item
  Deduce that \(\operatorname{Trd}_{\mathrm{D}/\mathrm{K}}(x)\) (resp. \(\operatorname{Nrd}_{\mathrm{D}/\mathrm{K}}(x)\) ) is the sum (resp. product) of m elements of the form \(txt^{-1}\) (apply Proposition 8 of VIII, p.~347).
\end{enumerate}

\begin{enumerate}
\def\labelenumi{\arabic{enumi})}
\setcounter{enumi}{6}
\tightlist
\item
  Let D be a field with center K of finite rank over K, endowed with a total order compatible with its additive group structure and such that the product of two positive elements is positive. Prove that D is equal to K. (Prove that K has characteristic zero. If \(D \neq K\) , observe that there exist elements x in D-K such that \(\operatorname{Trd}_{\mathrm{D}/\mathrm{K}}(x) = 0\) , and apply Exercise 6, d)).
\end{enumerate}

